% table10.tex -- provenance of Table~\ref{tab:native_phases} (Table 10).
%
% This is a provenance note, not body text. It records what produced the
% DeepSeek numbers in Table 10, traced 2026-11-03 from the sealed receipts.
% Include with % table10.tex -- provenance of Table~\ref{tab:native_phases} (Table 10).
%
% This is a provenance note, not body text. It records what produced the
% DeepSeek numbers in Table 10, traced 2026-11-03 from the sealed receipts.
% Include with % table10.tex -- provenance of Table~\ref{tab:native_phases} (Table 10).
%
% This is a provenance note, not body text. It records what produced the
% DeepSeek numbers in Table 10, traced 2026-11-03 from the sealed receipts.
% Include with % table10.tex -- provenance of Table~\ref{tab:native_phases} (Table 10).
%
% This is a provenance note, not body text. It records what produced the
% DeepSeek numbers in Table 10, traced 2026-11-03 from the sealed receipts.
% Include with \input{table10} inside the appendix if wanted; it is written as
% a self-contained section and is not referenced by any other file.
%
% Collected receipts: /media/miria/Elements/table10_provenance_20261103/
% (INDEX.md, MANIFEST.json, and the execution/comparison/fixed_chip_LP receipts).

\section{Provenance of the DeepSeek phase table}
\label{app:table10_provenance}

Table~\ref{tab:native_phases} is Table~10 of this paper, confirmed from
\path{neurips_2026.aux}. Its source is \path{update/decode_decomp.tex}. This
section records what produced its numbers, because the answer is not what the
surrounding text implies.

\paragraph{No LP solve produced the table.}
The custom-chip column is a \emph{replayed saved retiming} of an
already-admitted execution, evaluated by the analytical service evaluator. The
sealed closure receipt records that row under the key
\path{all_results.old_uniform_retiming} and reports
\path{new_optimizer_calls}~$=0$. The producing audit describes it as
``replayed saved uniform retiming'' and states that no optimizer, hardware
search, numerical format or kernel implementation was used. A fixed-chip LP
certificate does exist for the same chip, and is cited one paragraph below the
table, but it is a different quantity: its continuous upper lies a factor
$2.7055$ \emph{below} the latency printed in the table, it was also obtained
without an optimizer call, and the chip itself lies outside the certified LP
family.

\subsection{The three execution variants}
The closure receipt carries all three analytical executions of the same task
(32 requests, 4,096 prompt and 512 output tokens, 61 layers, 61,985,244 logical
invocations). Only the first is in Table~\ref{tab:native_phases}.

\begin{table}[t]
\centering\small
\caption{Analytical executions in the sealed closure receipt (ms). The first row
is Table~\ref{tab:native_phases}; the second supersedes it; the third is the
only genuinely LP-derived hardware.}
\label{tab:table10_variants}
\begin{tabular}{@{}lrrrr@{}}
\toprule
Receipt key & Prefill & Decode & Conversion & Total \\
\midrule
\path{old_uniform_retiming} & 1,078,913.842 & 2,074,998.907 & 42,972.849 & 3,196,885.598 \\
\path{retained_integral} & 913,989.589 & 2,074,998.907 & 39,920.585 & \textbf{3,028,909.081} \\
\path{rounded_LP_hardware} & 1,047,788.641 & 3,405,889.963 & 40,747.951 & 4,494,426.555 \\
\bottomrule
\end{tabular}
\end{table}

\paragraph{The table is one revision stale.}
The same 2026-10-27 closure replaced the unnecessary uniform dilation of prefill
and conversion kernel services with the supported scalar evaluator, tightening
the bound by $167{,}976.516589$~ms on the same design. Decode was untouched
(\path{decode_unchanged}~$=$~true), which is why Table~\ref{tab:native_phases}
still agrees on decode and no longer agrees on prefill or preparation. Against
the same reference column the two give

\begin{table}[t]
\centering\small
\caption{Printed versus current ratios against the identical reference column.}
\label{tab:table10_stale}
\begin{tabular}{@{}lrr@{}}
\toprule
Phase & Table~\ref{tab:native_phases} & Current receipt \\
\midrule
Prefill & $12.408\times$ & $\mathbf{14.647\times}$ \\
Decode & $25.582\times$ & $25.582\times$ \\
Preparation & $3.054\times$ & $\mathbf{3.288\times}$ \\
\midrule
Total & $20.833\times$ & $\mathbf{21.989\times}$ \\
\bottomrule
\end{tabular}
\end{table}

\noindent
so the printed table understates its own result. The LP appendix also still
cites the superseded $3{,}196{,}885.598$-ms latency.

\paragraph{The LP-derived hardware is slower.}
The rounded-LP candidate uses 128 arrays, 4,096 vector lanes, 1,597~GB/s DRAM
and 1,024~GB/s host, and passes the same physical scenario at $313.651658$~mm$^2$
and $298.3871$~W. Its execution is $4{,}494{,}426.555$~ms, a factor $1.4838$
slower than the retained chip; its untightened bound was $5{,}509{,}427.509$~ms.
The larger recovered chip was retained on that basis.

\subsection{The reference column and its sensitivity}
The reference column of Table~\ref{tab:native_phases} is the retained
host-backed analytical extension: $13{,}387{,}192.863$~ms prefill,
$53{,}083{,}323.571$~ms decode, $131{,}252.321$~ms conversion,
$66{,}601{,}768.756$~ms total. Its comparison receipt carries
\path{numerical_optimizer_calls}~$=0$, \path{inference_calls}~$=0$,
\path{headline_speedup_admitted}~$=$~false,
\path{verified_FAST_modeled_TPUv3_comparison}~$=$~false and
\path{FAST_method_comparison}~$=$~false.

The headline depends on an assumed reference interface rate. The receipt's
sensitivity set holds eighteen cases, nine at each assumed host rate; recomputed
from its exact fractions against the current $3{,}028{,}909.081$-ms execution:

\begin{table}[t]
\centering\small
\caption{Overall ratio under the receipt's own reference sensitivity set.}
\label{tab:table10_sensitivity}
\begin{tabular}{@{}lrr@{}}
\toprule
Assumed reference host rate & Range & Median \\
\midrule
32~GB/s & $21.90491$--$22.15628\times$ & $21.98870\times$ \\
1,024~GB/s & $4.42769$--$4.67907\times$ & $4.51148\times$ \\
\bottomrule
\end{tabular}
\end{table}

\noindent
Table~\ref{tab:native_phases} uses the 32~GB/s case with 8,192~B/ns local
service and 2,048 vector primitive cycles/ns. The result therefore moves by a
factor of roughly $4.9$ on a reference interface assumption alone, which is why
the producing audit declines to promote it to a FAST-baseline headline.

\subsection{The fixed-chip LP, and why it is a separate fact}
The certificate cited beside the table has source identity
\path{577942e1db9ffbe4f6582d27a69a418765e60deec584e04465f70b4f7a558604} over the
original source
\path{831edc6f8cde9bd2091541b55beaa00fae2736035badb3937236ba32c66104e3}. It
retains all $15{,}044{,}386{,}833$ implementation choices, $61{,}985{,}244$
invocations and $5{,}809{,}955{,}157$ original event labels as independent, and
fixes the hardware to the chip in Table~\ref{tab:native_phases}. Its checked
endpoints are
\[
 \begin{gathered}
 L=1{,}180{,}606.049320422\ \mathrm{ms},\qquad
 U=1{,}181{,}606.6831816933\ \mathrm{ms},\\
 100(U-L)/U=0.0846841741\%,
 \end{gathered}
\]
with status \path{COMPLETE_FIXED_CHIP_NATIVE_INTERVAL} and
\path{optimizer_calls}~$=0$. The lower is constructed analytically as
$L=\kappa\sum_i\pi_i/B_h$ with $B_h=1630$, from the exact resource row
$Q^B_e\le B_h d_e$ and the already-checked coverage prices $\pi_i$; it transfers
a row proof rather than an earlier scalar bound. The upper retains the admitted
point's candidate, lifecycle, storage and translation fields, recomputes
services at the new rates, and takes the resource maximum after temporal mixing
under a conservative host factor $8/5$. The checker covers $197{,}075$ source
service cases, $791{,}997$ shared row instances and all $19{,}317$ point service
classes; admission takes $94.626$~s and the replay's admission $96.503$~s.

Three facts keep this certificate separate from Table~\ref{tab:native_phases}.

\begin{enumerate}
\item \textbf{Scale.} $U=1{,}181.607$~s against the table's $3{,}196.886$~s, a
factor $2.7055$. The continuous upper is not an executable latency.
\item \textbf{The chip is outside the certified family.} The retained chip uses
20 host interface units, while the frozen target
\path{1946a0a176e74c6fbad85a88b58d83e95adf6beeaecf98bfa1519866acad91c0} contains
fixed host-interface branches with 1, 32 or 34 units only. The LP was not
enlarged, so no executable optimality interval follows from pairing its lower
with this execution.
\item \textbf{The LP's own hardware loses.} See
Table~\ref{tab:table10_variants}, third row.
\end{enumerate}

\noindent
Recorded alongside: the relaxed fractional witness is $1{,}205{,}554.982396$~ms,
and at exact installed rates the same fractional choices cost
$1{,}477{,}032.945143$~ms. Neither fractional point is the integral execution.

\subsection{What actually selected the chip and schedule}
A recovery-and-replay pipeline, not a solve.

\begin{enumerate}
\item An integral chip and complete execution were recovered earlier and saved.
\item The 2026-10-27 closure replayed that saved design against the current
physical screens: $61{,}985{,}244$ logical invocations, $68{,}800{,}988$ original
owners, $13{,}042{,}436$ weighted occurrences, $281{,}872$ retained gaps,
$19{,}230$ original DRAM imports and leases, and all $384$ selected pipeline
recipe types regenerated and independently replayed. Capacity was recounted:
decode DRAM $68{,}462{,}901{,}448$ of $68{,}719{,}476{,}736$ bytes, SRAM
$62{,}078{,}976$ of $82{,}051{,}072$, host address space
$2{,}046{,}265{,}271{,}296$ of $2{,}199{,}023{,}255{,}552$.
\item Phase accounting was then tightened as described above.
\item The fixed-chip LP interval was constructed separately for the same chip.
\end{enumerate}

\noindent
The chip has 320 $16\times16$ BF16 MAC arrays, 5,248 vector lanes, 1,024
conversion lanes, 78.25~MiB SRAM, 64~GiB DRAM and 2~TiB host backing, with
1,630~GB/s DRAM, 640~GB/s host and 20,480~GB/s local service at a 1~ns clock. It
evaluates to $321.686108$~mm$^2$ and $298.6865$~W under the BF16 area-scale-four
scenario and passes every physical predicate there. These remain uncalibrated
architectural coefficients, not physical-design sign-off.

\subsection{How to quote the table}
\begin{itemize}
\item The decode ratio $25.582\times$ is current and survives the retiming
change.
\item The total should be $21.989\times$ and prefill $14.647\times$; the printed
$20.833\times$ and $12.408\times$ predate the 2026-10-27 tightening.
\item Both are measured against an assumed 32~GB/s host reference. At
1,024~GB/s the same design gives $4.511\times$. The assumption belongs in the
same sentence as the ratio.
\item It is not an LP optimum, not a certified gap and not a comparison against
FAST's modeled TPU-v3 baseline; the receipt denies the last two explicitly.
\end{itemize}
 inside the appendix if wanted; it is written as
% a self-contained section and is not referenced by any other file.
%
% Collected receipts: /media/miria/Elements/table10_provenance_20261103/
% (INDEX.md, MANIFEST.json, and the execution/comparison/fixed_chip_LP receipts).

\section{Provenance of the DeepSeek phase table}
\label{app:table10_provenance}

Table~\ref{tab:native_phases} is Table~10 of this paper, confirmed from
\path{neurips_2026.aux}. Its source is \path{update/decode_decomp.tex}. This
section records what produced its numbers, because the answer is not what the
surrounding text implies.

\paragraph{No LP solve produced the table.}
The custom-chip column is a \emph{replayed saved retiming} of an
already-admitted execution, evaluated by the analytical service evaluator. The
sealed closure receipt records that row under the key
\path{all_results.old_uniform_retiming} and reports
\path{new_optimizer_calls}~$=0$. The producing audit describes it as
``replayed saved uniform retiming'' and states that no optimizer, hardware
search, numerical format or kernel implementation was used. A fixed-chip LP
certificate does exist for the same chip, and is cited one paragraph below the
table, but it is a different quantity: its continuous upper lies a factor
$2.7055$ \emph{below} the latency printed in the table, it was also obtained
without an optimizer call, and the chip itself lies outside the certified LP
family.

\subsection{The three execution variants}
The closure receipt carries all three analytical executions of the same task
(32 requests, 4,096 prompt and 512 output tokens, 61 layers, 61,985,244 logical
invocations). Only the first is in Table~\ref{tab:native_phases}.

\begin{table}[t]
\centering\small
\caption{Analytical executions in the sealed closure receipt (ms). The first row
is Table~\ref{tab:native_phases}; the second supersedes it; the third is the
only genuinely LP-derived hardware.}
\label{tab:table10_variants}
\begin{tabular}{@{}lrrrr@{}}
\toprule
Receipt key & Prefill & Decode & Conversion & Total \\
\midrule
\path{old_uniform_retiming} & 1,078,913.842 & 2,074,998.907 & 42,972.849 & 3,196,885.598 \\
\path{retained_integral} & 913,989.589 & 2,074,998.907 & 39,920.585 & \textbf{3,028,909.081} \\
\path{rounded_LP_hardware} & 1,047,788.641 & 3,405,889.963 & 40,747.951 & 4,494,426.555 \\
\bottomrule
\end{tabular}
\end{table}

\paragraph{The table is one revision stale.}
The same 2026-10-27 closure replaced the unnecessary uniform dilation of prefill
and conversion kernel services with the supported scalar evaluator, tightening
the bound by $167{,}976.516589$~ms on the same design. Decode was untouched
(\path{decode_unchanged}~$=$~true), which is why Table~\ref{tab:native_phases}
still agrees on decode and no longer agrees on prefill or preparation. Against
the same reference column the two give

\begin{table}[t]
\centering\small
\caption{Printed versus current ratios against the identical reference column.}
\label{tab:table10_stale}
\begin{tabular}{@{}lrr@{}}
\toprule
Phase & Table~\ref{tab:native_phases} & Current receipt \\
\midrule
Prefill & $12.408\times$ & $\mathbf{14.647\times}$ \\
Decode & $25.582\times$ & $25.582\times$ \\
Preparation & $3.054\times$ & $\mathbf{3.288\times}$ \\
\midrule
Total & $20.833\times$ & $\mathbf{21.989\times}$ \\
\bottomrule
\end{tabular}
\end{table}

\noindent
so the printed table understates its own result. The LP appendix also still
cites the superseded $3{,}196{,}885.598$-ms latency.

\paragraph{The LP-derived hardware is slower.}
The rounded-LP candidate uses 128 arrays, 4,096 vector lanes, 1,597~GB/s DRAM
and 1,024~GB/s host, and passes the same physical scenario at $313.651658$~mm$^2$
and $298.3871$~W. Its execution is $4{,}494{,}426.555$~ms, a factor $1.4838$
slower than the retained chip; its untightened bound was $5{,}509{,}427.509$~ms.
The larger recovered chip was retained on that basis.

\subsection{The reference column and its sensitivity}
The reference column of Table~\ref{tab:native_phases} is the retained
host-backed analytical extension: $13{,}387{,}192.863$~ms prefill,
$53{,}083{,}323.571$~ms decode, $131{,}252.321$~ms conversion,
$66{,}601{,}768.756$~ms total. Its comparison receipt carries
\path{numerical_optimizer_calls}~$=0$, \path{inference_calls}~$=0$,
\path{headline_speedup_admitted}~$=$~false,
\path{verified_FAST_modeled_TPUv3_comparison}~$=$~false and
\path{FAST_method_comparison}~$=$~false.

The headline depends on an assumed reference interface rate. The receipt's
sensitivity set holds eighteen cases, nine at each assumed host rate; recomputed
from its exact fractions against the current $3{,}028{,}909.081$-ms execution:

\begin{table}[t]
\centering\small
\caption{Overall ratio under the receipt's own reference sensitivity set.}
\label{tab:table10_sensitivity}
\begin{tabular}{@{}lrr@{}}
\toprule
Assumed reference host rate & Range & Median \\
\midrule
32~GB/s & $21.90491$--$22.15628\times$ & $21.98870\times$ \\
1,024~GB/s & $4.42769$--$4.67907\times$ & $4.51148\times$ \\
\bottomrule
\end{tabular}
\end{table}

\noindent
Table~\ref{tab:native_phases} uses the 32~GB/s case with 8,192~B/ns local
service and 2,048 vector primitive cycles/ns. The result therefore moves by a
factor of roughly $4.9$ on a reference interface assumption alone, which is why
the producing audit declines to promote it to a FAST-baseline headline.

\subsection{The fixed-chip LP, and why it is a separate fact}
The certificate cited beside the table has source identity
\path{577942e1db9ffbe4f6582d27a69a418765e60deec584e04465f70b4f7a558604} over the
original source
\path{831edc6f8cde9bd2091541b55beaa00fae2736035badb3937236ba32c66104e3}. It
retains all $15{,}044{,}386{,}833$ implementation choices, $61{,}985{,}244$
invocations and $5{,}809{,}955{,}157$ original event labels as independent, and
fixes the hardware to the chip in Table~\ref{tab:native_phases}. Its checked
endpoints are
\[
 \begin{gathered}
 L=1{,}180{,}606.049320422\ \mathrm{ms},\qquad
 U=1{,}181{,}606.6831816933\ \mathrm{ms},\\
 100(U-L)/U=0.0846841741\%,
 \end{gathered}
\]
with status \path{COMPLETE_FIXED_CHIP_NATIVE_INTERVAL} and
\path{optimizer_calls}~$=0$. The lower is constructed analytically as
$L=\kappa\sum_i\pi_i/B_h$ with $B_h=1630$, from the exact resource row
$Q^B_e\le B_h d_e$ and the already-checked coverage prices $\pi_i$; it transfers
a row proof rather than an earlier scalar bound. The upper retains the admitted
point's candidate, lifecycle, storage and translation fields, recomputes
services at the new rates, and takes the resource maximum after temporal mixing
under a conservative host factor $8/5$. The checker covers $197{,}075$ source
service cases, $791{,}997$ shared row instances and all $19{,}317$ point service
classes; admission takes $94.626$~s and the replay's admission $96.503$~s.

Three facts keep this certificate separate from Table~\ref{tab:native_phases}.

\begin{enumerate}
\item \textbf{Scale.} $U=1{,}181.607$~s against the table's $3{,}196.886$~s, a
factor $2.7055$. The continuous upper is not an executable latency.
\item \textbf{The chip is outside the certified family.} The retained chip uses
20 host interface units, while the frozen target
\path{1946a0a176e74c6fbad85a88b58d83e95adf6beeaecf98bfa1519866acad91c0} contains
fixed host-interface branches with 1, 32 or 34 units only. The LP was not
enlarged, so no executable optimality interval follows from pairing its lower
with this execution.
\item \textbf{The LP's own hardware loses.} See
Table~\ref{tab:table10_variants}, third row.
\end{enumerate}

\noindent
Recorded alongside: the relaxed fractional witness is $1{,}205{,}554.982396$~ms,
and at exact installed rates the same fractional choices cost
$1{,}477{,}032.945143$~ms. Neither fractional point is the integral execution.

\subsection{What actually selected the chip and schedule}
A recovery-and-replay pipeline, not a solve.

\begin{enumerate}
\item An integral chip and complete execution were recovered earlier and saved.
\item The 2026-10-27 closure replayed that saved design against the current
physical screens: $61{,}985{,}244$ logical invocations, $68{,}800{,}988$ original
owners, $13{,}042{,}436$ weighted occurrences, $281{,}872$ retained gaps,
$19{,}230$ original DRAM imports and leases, and all $384$ selected pipeline
recipe types regenerated and independently replayed. Capacity was recounted:
decode DRAM $68{,}462{,}901{,}448$ of $68{,}719{,}476{,}736$ bytes, SRAM
$62{,}078{,}976$ of $82{,}051{,}072$, host address space
$2{,}046{,}265{,}271{,}296$ of $2{,}199{,}023{,}255{,}552$.
\item Phase accounting was then tightened as described above.
\item The fixed-chip LP interval was constructed separately for the same chip.
\end{enumerate}

\noindent
The chip has 320 $16\times16$ BF16 MAC arrays, 5,248 vector lanes, 1,024
conversion lanes, 78.25~MiB SRAM, 64~GiB DRAM and 2~TiB host backing, with
1,630~GB/s DRAM, 640~GB/s host and 20,480~GB/s local service at a 1~ns clock. It
evaluates to $321.686108$~mm$^2$ and $298.6865$~W under the BF16 area-scale-four
scenario and passes every physical predicate there. These remain uncalibrated
architectural coefficients, not physical-design sign-off.

\subsection{How to quote the table}
\begin{itemize}
\item The decode ratio $25.582\times$ is current and survives the retiming
change.
\item The total should be $21.989\times$ and prefill $14.647\times$; the printed
$20.833\times$ and $12.408\times$ predate the 2026-10-27 tightening.
\item Both are measured against an assumed 32~GB/s host reference. At
1,024~GB/s the same design gives $4.511\times$. The assumption belongs in the
same sentence as the ratio.
\item It is not an LP optimum, not a certified gap and not a comparison against
FAST's modeled TPU-v3 baseline; the receipt denies the last two explicitly.
\end{itemize}
 inside the appendix if wanted; it is written as
% a self-contained section and is not referenced by any other file.
%
% Collected receipts: /media/miria/Elements/table10_provenance_20261103/
% (INDEX.md, MANIFEST.json, and the execution/comparison/fixed_chip_LP receipts).

\section{Provenance of the DeepSeek phase table}
\label{app:table10_provenance}

Table~\ref{tab:native_phases} is Table~10 of this paper, confirmed from
\path{neurips_2026.aux}. Its source is \path{update/decode_decomp.tex}. This
section records what produced its numbers, because the answer is not what the
surrounding text implies.

\paragraph{No LP solve produced the table.}
The custom-chip column is a \emph{replayed saved retiming} of an
already-admitted execution, evaluated by the analytical service evaluator. The
sealed closure receipt records that row under the key
\path{all_results.old_uniform_retiming} and reports
\path{new_optimizer_calls}~$=0$. The producing audit describes it as
``replayed saved uniform retiming'' and states that no optimizer, hardware
search, numerical format or kernel implementation was used. A fixed-chip LP
certificate does exist for the same chip, and is cited one paragraph below the
table, but it is a different quantity: its continuous upper lies a factor
$2.7055$ \emph{below} the latency printed in the table, it was also obtained
without an optimizer call, and the chip itself lies outside the certified LP
family.

\subsection{The three execution variants}
The closure receipt carries all three analytical executions of the same task
(32 requests, 4,096 prompt and 512 output tokens, 61 layers, 61,985,244 logical
invocations). Only the first is in Table~\ref{tab:native_phases}.

\begin{table}[t]
\centering\small
\caption{Analytical executions in the sealed closure receipt (ms). The first row
is Table~\ref{tab:native_phases}; the second supersedes it; the third is the
only genuinely LP-derived hardware.}
\label{tab:table10_variants}
\begin{tabular}{@{}lrrrr@{}}
\toprule
Receipt key & Prefill & Decode & Conversion & Total \\
\midrule
\path{old_uniform_retiming} & 1,078,913.842 & 2,074,998.907 & 42,972.849 & 3,196,885.598 \\
\path{retained_integral} & 913,989.589 & 2,074,998.907 & 39,920.585 & \textbf{3,028,909.081} \\
\path{rounded_LP_hardware} & 1,047,788.641 & 3,405,889.963 & 40,747.951 & 4,494,426.555 \\
\bottomrule
\end{tabular}
\end{table}

\paragraph{The table is one revision stale.}
The same 2026-10-27 closure replaced the unnecessary uniform dilation of prefill
and conversion kernel services with the supported scalar evaluator, tightening
the bound by $167{,}976.516589$~ms on the same design. Decode was untouched
(\path{decode_unchanged}~$=$~true), which is why Table~\ref{tab:native_phases}
still agrees on decode and no longer agrees on prefill or preparation. Against
the same reference column the two give

\begin{table}[t]
\centering\small
\caption{Printed versus current ratios against the identical reference column.}
\label{tab:table10_stale}
\begin{tabular}{@{}lrr@{}}
\toprule
Phase & Table~\ref{tab:native_phases} & Current receipt \\
\midrule
Prefill & $12.408\times$ & $\mathbf{14.647\times}$ \\
Decode & $25.582\times$ & $25.582\times$ \\
Preparation & $3.054\times$ & $\mathbf{3.288\times}$ \\
\midrule
Total & $20.833\times$ & $\mathbf{21.989\times}$ \\
\bottomrule
\end{tabular}
\end{table}

\noindent
so the printed table understates its own result. The LP appendix also still
cites the superseded $3{,}196{,}885.598$-ms latency.

\paragraph{The LP-derived hardware is slower.}
The rounded-LP candidate uses 128 arrays, 4,096 vector lanes, 1,597~GB/s DRAM
and 1,024~GB/s host, and passes the same physical scenario at $313.651658$~mm$^2$
and $298.3871$~W. Its execution is $4{,}494{,}426.555$~ms, a factor $1.4838$
slower than the retained chip; its untightened bound was $5{,}509{,}427.509$~ms.
The larger recovered chip was retained on that basis.

\subsection{The reference column and its sensitivity}
The reference column of Table~\ref{tab:native_phases} is the retained
host-backed analytical extension: $13{,}387{,}192.863$~ms prefill,
$53{,}083{,}323.571$~ms decode, $131{,}252.321$~ms conversion,
$66{,}601{,}768.756$~ms total. Its comparison receipt carries
\path{numerical_optimizer_calls}~$=0$, \path{inference_calls}~$=0$,
\path{headline_speedup_admitted}~$=$~false,
\path{verified_FAST_modeled_TPUv3_comparison}~$=$~false and
\path{FAST_method_comparison}~$=$~false.

The headline depends on an assumed reference interface rate. The receipt's
sensitivity set holds eighteen cases, nine at each assumed host rate; recomputed
from its exact fractions against the current $3{,}028{,}909.081$-ms execution:

\begin{table}[t]
\centering\small
\caption{Overall ratio under the receipt's own reference sensitivity set.}
\label{tab:table10_sensitivity}
\begin{tabular}{@{}lrr@{}}
\toprule
Assumed reference host rate & Range & Median \\
\midrule
32~GB/s & $21.90491$--$22.15628\times$ & $21.98870\times$ \\
1,024~GB/s & $4.42769$--$4.67907\times$ & $4.51148\times$ \\
\bottomrule
\end{tabular}
\end{table}

\noindent
Table~\ref{tab:native_phases} uses the 32~GB/s case with 8,192~B/ns local
service and 2,048 vector primitive cycles/ns. The result therefore moves by a
factor of roughly $4.9$ on a reference interface assumption alone, which is why
the producing audit declines to promote it to a FAST-baseline headline.

\subsection{The fixed-chip LP, and why it is a separate fact}
The certificate cited beside the table has source identity
\path{577942e1db9ffbe4f6582d27a69a418765e60deec584e04465f70b4f7a558604} over the
original source
\path{831edc6f8cde9bd2091541b55beaa00fae2736035badb3937236ba32c66104e3}. It
retains all $15{,}044{,}386{,}833$ implementation choices, $61{,}985{,}244$
invocations and $5{,}809{,}955{,}157$ original event labels as independent, and
fixes the hardware to the chip in Table~\ref{tab:native_phases}. Its checked
endpoints are
\[
 \begin{gathered}
 L=1{,}180{,}606.049320422\ \mathrm{ms},\qquad
 U=1{,}181{,}606.6831816933\ \mathrm{ms},\\
 100(U-L)/U=0.0846841741\%,
 \end{gathered}
\]
with status \path{COMPLETE_FIXED_CHIP_NATIVE_INTERVAL} and
\path{optimizer_calls}~$=0$. The lower is constructed analytically as
$L=\kappa\sum_i\pi_i/B_h$ with $B_h=1630$, from the exact resource row
$Q^B_e\le B_h d_e$ and the already-checked coverage prices $\pi_i$; it transfers
a row proof rather than an earlier scalar bound. The upper retains the admitted
point's candidate, lifecycle, storage and translation fields, recomputes
services at the new rates, and takes the resource maximum after temporal mixing
under a conservative host factor $8/5$. The checker covers $197{,}075$ source
service cases, $791{,}997$ shared row instances and all $19{,}317$ point service
classes; admission takes $94.626$~s and the replay's admission $96.503$~s.

Three facts keep this certificate separate from Table~\ref{tab:native_phases}.

\begin{enumerate}
\item \textbf{Scale.} $U=1{,}181.607$~s against the table's $3{,}196.886$~s, a
factor $2.7055$. The continuous upper is not an executable latency.
\item \textbf{The chip is outside the certified family.} The retained chip uses
20 host interface units, while the frozen target
\path{1946a0a176e74c6fbad85a88b58d83e95adf6beeaecf98bfa1519866acad91c0} contains
fixed host-interface branches with 1, 32 or 34 units only. The LP was not
enlarged, so no executable optimality interval follows from pairing its lower
with this execution.
\item \textbf{The LP's own hardware loses.} See
Table~\ref{tab:table10_variants}, third row.
\end{enumerate}

\noindent
Recorded alongside: the relaxed fractional witness is $1{,}205{,}554.982396$~ms,
and at exact installed rates the same fractional choices cost
$1{,}477{,}032.945143$~ms. Neither fractional point is the integral execution.

\subsection{What actually selected the chip and schedule}
A recovery-and-replay pipeline, not a solve.

\begin{enumerate}
\item An integral chip and complete execution were recovered earlier and saved.
\item The 2026-10-27 closure replayed that saved design against the current
physical screens: $61{,}985{,}244$ logical invocations, $68{,}800{,}988$ original
owners, $13{,}042{,}436$ weighted occurrences, $281{,}872$ retained gaps,
$19{,}230$ original DRAM imports and leases, and all $384$ selected pipeline
recipe types regenerated and independently replayed. Capacity was recounted:
decode DRAM $68{,}462{,}901{,}448$ of $68{,}719{,}476{,}736$ bytes, SRAM
$62{,}078{,}976$ of $82{,}051{,}072$, host address space
$2{,}046{,}265{,}271{,}296$ of $2{,}199{,}023{,}255{,}552$.
\item Phase accounting was then tightened as described above.
\item The fixed-chip LP interval was constructed separately for the same chip.
\end{enumerate}

\noindent
The chip has 320 $16\times16$ BF16 MAC arrays, 5,248 vector lanes, 1,024
conversion lanes, 78.25~MiB SRAM, 64~GiB DRAM and 2~TiB host backing, with
1,630~GB/s DRAM, 640~GB/s host and 20,480~GB/s local service at a 1~ns clock. It
evaluates to $321.686108$~mm$^2$ and $298.6865$~W under the BF16 area-scale-four
scenario and passes every physical predicate there. These remain uncalibrated
architectural coefficients, not physical-design sign-off.

\subsection{How to quote the table}
\begin{itemize}
\item The decode ratio $25.582\times$ is current and survives the retiming
change.
\item The total should be $21.989\times$ and prefill $14.647\times$; the printed
$20.833\times$ and $12.408\times$ predate the 2026-10-27 tightening.
\item Both are measured against an assumed 32~GB/s host reference. At
1,024~GB/s the same design gives $4.511\times$. The assumption belongs in the
same sentence as the ratio.
\item It is not an LP optimum, not a certified gap and not a comparison against
FAST's modeled TPU-v3 baseline; the receipt denies the last two explicitly.
\end{itemize}
 inside the appendix if wanted; it is written as
% a self-contained section and is not referenced by any other file.
%
% Collected receipts: /media/miria/Elements/table10_provenance_20261103/
% (INDEX.md, MANIFEST.json, and the execution/comparison/fixed_chip_LP receipts).

\section{Provenance of the DeepSeek phase table}
\label{app:table10_provenance}

Table~\ref{tab:native_phases} is Table~10 of this paper, confirmed from
\path{neurips_2026.aux}. Its source is \path{update/decode_decomp.tex}. This
section records what produced its numbers, because the answer is not what the
surrounding text implies.

\paragraph{No LP solve produced the table.}
The custom-chip column is a \emph{replayed saved retiming} of an
already-admitted execution, evaluated by the analytical service evaluator. The
sealed closure receipt records that row under the key
\path{all_results.old_uniform_retiming} and reports
\path{new_optimizer_calls}~$=0$. The producing audit describes it as
``replayed saved uniform retiming'' and states that no optimizer, hardware
search, numerical format or kernel implementation was used. A fixed-chip LP
certificate does exist for the same chip, and is cited one paragraph below the
table, but it is a different quantity: its continuous upper lies a factor
$2.7055$ \emph{below} the latency printed in the table, it was also obtained
without an optimizer call, and the chip itself lies outside the certified LP
family.

\subsection{The three execution variants}
The closure receipt carries all three analytical executions of the same task
(32 requests, 4,096 prompt and 512 output tokens, 61 layers, 61,985,244 logical
invocations). Only the first is in Table~\ref{tab:native_phases}.

\begin{table}[t]
\centering\small
\caption{Analytical executions in the sealed closure receipt (ms). The first row
is Table~\ref{tab:native_phases}; the second supersedes it; the third is the
only genuinely LP-derived hardware.}
\label{tab:table10_variants}
\begin{tabular}{@{}lrrrr@{}}
\toprule
Receipt key & Prefill & Decode & Conversion & Total \\
\midrule
\path{old_uniform_retiming} & 1,078,913.842 & 2,074,998.907 & 42,972.849 & 3,196,885.598 \\
\path{retained_integral} & 913,989.589 & 2,074,998.907 & 39,920.585 & \textbf{3,028,909.081} \\
\path{rounded_LP_hardware} & 1,047,788.641 & 3,405,889.963 & 40,747.951 & 4,494,426.555 \\
\bottomrule
\end{tabular}
\end{table}

\paragraph{The table is one revision stale.}
The same 2026-10-27 closure replaced the unnecessary uniform dilation of prefill
and conversion kernel services with the supported scalar evaluator, tightening
the bound by $167{,}976.516589$~ms on the same design. Decode was untouched
(\path{decode_unchanged}~$=$~true), which is why Table~\ref{tab:native_phases}
still agrees on decode and no longer agrees on prefill or preparation. Against
the same reference column the two give

\begin{table}[t]
\centering\small
\caption{Printed versus current ratios against the identical reference column.}
\label{tab:table10_stale}
\begin{tabular}{@{}lrr@{}}
\toprule
Phase & Table~\ref{tab:native_phases} & Current receipt \\
\midrule
Prefill & $12.408\times$ & $\mathbf{14.647\times}$ \\
Decode & $25.582\times$ & $25.582\times$ \\
Preparation & $3.054\times$ & $\mathbf{3.288\times}$ \\
\midrule
Total & $20.833\times$ & $\mathbf{21.989\times}$ \\
\bottomrule
\end{tabular}
\end{table}

\noindent
so the printed table understates its own result. The LP appendix also still
cites the superseded $3{,}196{,}885.598$-ms latency.

\paragraph{The LP-derived hardware is slower.}
The rounded-LP candidate uses 128 arrays, 4,096 vector lanes, 1,597~GB/s DRAM
and 1,024~GB/s host, and passes the same physical scenario at $313.651658$~mm$^2$
and $298.3871$~W. Its execution is $4{,}494{,}426.555$~ms, a factor $1.4838$
slower than the retained chip; its untightened bound was $5{,}509{,}427.509$~ms.
The larger recovered chip was retained on that basis.

\subsection{The reference column and its sensitivity}
The reference column of Table~\ref{tab:native_phases} is the retained
host-backed analytical extension: $13{,}387{,}192.863$~ms prefill,
$53{,}083{,}323.571$~ms decode, $131{,}252.321$~ms conversion,
$66{,}601{,}768.756$~ms total. Its comparison receipt carries
\path{numerical_optimizer_calls}~$=0$, \path{inference_calls}~$=0$,
\path{headline_speedup_admitted}~$=$~false,
\path{verified_FAST_modeled_TPUv3_comparison}~$=$~false and
\path{FAST_method_comparison}~$=$~false.

The headline depends on an assumed reference interface rate. The receipt's
sensitivity set holds eighteen cases, nine at each assumed host rate; recomputed
from its exact fractions against the current $3{,}028{,}909.081$-ms execution:

\begin{table}[t]
\centering\small
\caption{Overall ratio under the receipt's own reference sensitivity set.}
\label{tab:table10_sensitivity}
\begin{tabular}{@{}lrr@{}}
\toprule
Assumed reference host rate & Range & Median \\
\midrule
32~GB/s & $21.90491$--$22.15628\times$ & $21.98870\times$ \\
1,024~GB/s & $4.42769$--$4.67907\times$ & $4.51148\times$ \\
\bottomrule
\end{tabular}
\end{table}

\noindent
Table~\ref{tab:native_phases} uses the 32~GB/s case with 8,192~B/ns local
service and 2,048 vector primitive cycles/ns. The result therefore moves by a
factor of roughly $4.9$ on a reference interface assumption alone, which is why
the producing audit declines to promote it to a FAST-baseline headline.

\subsection{The fixed-chip LP, and why it is a separate fact}
The certificate cited beside the table has source identity
\path{577942e1db9ffbe4f6582d27a69a418765e60deec584e04465f70b4f7a558604} over the
original source
\path{831edc6f8cde9bd2091541b55beaa00fae2736035badb3937236ba32c66104e3}. It
retains all $15{,}044{,}386{,}833$ implementation choices, $61{,}985{,}244$
invocations and $5{,}809{,}955{,}157$ original event labels as independent, and
fixes the hardware to the chip in Table~\ref{tab:native_phases}. Its checked
endpoints are
\[
 \begin{gathered}
 L=1{,}180{,}606.049320422\ \mathrm{ms},\qquad
 U=1{,}181{,}606.6831816933\ \mathrm{ms},\\
 100(U-L)/U=0.0846841741\%,
 \end{gathered}
\]
with status \path{COMPLETE_FIXED_CHIP_NATIVE_INTERVAL} and
\path{optimizer_calls}~$=0$. The lower is constructed analytically as
$L=\kappa\sum_i\pi_i/B_h$ with $B_h=1630$, from the exact resource row
$Q^B_e\le B_h d_e$ and the already-checked coverage prices $\pi_i$; it transfers
a row proof rather than an earlier scalar bound. The upper retains the admitted
point's candidate, lifecycle, storage and translation fields, recomputes
services at the new rates, and takes the resource maximum after temporal mixing
under a conservative host factor $8/5$. The checker covers $197{,}075$ source
service cases, $791{,}997$ shared row instances and all $19{,}317$ point service
classes; admission takes $94.626$~s and the replay's admission $96.503$~s.

Three facts keep this certificate separate from Table~\ref{tab:native_phases}.

\begin{enumerate}
\item \textbf{Scale.} $U=1{,}181.607$~s against the table's $3{,}196.886$~s, a
factor $2.7055$. The continuous upper is not an executable latency.
\item \textbf{The chip is outside the certified family.} The retained chip uses
20 host interface units, while the frozen target
\path{1946a0a176e74c6fbad85a88b58d83e95adf6beeaecf98bfa1519866acad91c0} contains
fixed host-interface branches with 1, 32 or 34 units only. The LP was not
enlarged, so no executable optimality interval follows from pairing its lower
with this execution.
\item \textbf{The LP's own hardware loses.} See
Table~\ref{tab:table10_variants}, third row.
\end{enumerate}

\noindent
Recorded alongside: the relaxed fractional witness is $1{,}205{,}554.982396$~ms,
and at exact installed rates the same fractional choices cost
$1{,}477{,}032.945143$~ms. Neither fractional point is the integral execution.

\subsection{What actually selected the chip and schedule}
A recovery-and-replay pipeline, not a solve.

\begin{enumerate}
\item An integral chip and complete execution were recovered earlier and saved.
\item The 2026-10-27 closure replayed that saved design against the current
physical screens: $61{,}985{,}244$ logical invocations, $68{,}800{,}988$ original
owners, $13{,}042{,}436$ weighted occurrences, $281{,}872$ retained gaps,
$19{,}230$ original DRAM imports and leases, and all $384$ selected pipeline
recipe types regenerated and independently replayed. Capacity was recounted:
decode DRAM $68{,}462{,}901{,}448$ of $68{,}719{,}476{,}736$ bytes, SRAM
$62{,}078{,}976$ of $82{,}051{,}072$, host address space
$2{,}046{,}265{,}271{,}296$ of $2{,}199{,}023{,}255{,}552$.
\item Phase accounting was then tightened as described above.
\item The fixed-chip LP interval was constructed separately for the same chip.
\end{enumerate}

\noindent
The chip has 320 $16\times16$ BF16 MAC arrays, 5,248 vector lanes, 1,024
conversion lanes, 78.25~MiB SRAM, 64~GiB DRAM and 2~TiB host backing, with
1,630~GB/s DRAM, 640~GB/s host and 20,480~GB/s local service at a 1~ns clock. It
evaluates to $321.686108$~mm$^2$ and $298.6865$~W under the BF16 area-scale-four
scenario and passes every physical predicate there. These remain uncalibrated
architectural coefficients, not physical-design sign-off.

\subsection{How to quote the table}
\begin{itemize}
\item The decode ratio $25.582\times$ is current and survives the retiming
change.
\item The total should be $21.989\times$ and prefill $14.647\times$; the printed
$20.833\times$ and $12.408\times$ predate the 2026-10-27 tightening.
\item Both are measured against an assumed 32~GB/s host reference. At
1,024~GB/s the same design gives $4.511\times$. The assumption belongs in the
same sentence as the ratio.
\item It is not an LP optimum, not a certified gap and not a comparison against
FAST's modeled TPU-v3 baseline; the receipt denies the last two explicitly.
\end{itemize}
